%%%%%%%%%%  scan idx 289  |  folio 276  %%%%%%%%%%
induite par \(\cl_{K^b(A)/K^{b,\emptyset}(A)}\) (1.1) est une application additive universelle, i.e. que le couple \(\left(k\bigl(K^b(A)/K^{b,\emptyset}(A)\bigr),\cl_A\right)\) représente le foncteur associant à chaque groupe abélien \(G\) le groupe abélien des applications additives de \(\Ob A\) dans \(G\). On posera
\begin{equation}\tag{1.5.2}
 k\bigl(K^b(A)/K^{b,\emptyset}(A)\bigr)=k(A)
\end{equation}
(on écrira \(k(A,C)\) au lieu de \(k(A)\) quand on voudra éviter toute ambiguïté sur la catégorie abélienne ambiante)\footnote{La notation \(k(A)\) avait été introduite dans (I 6.3) pour désigner le groupe \(k(K^b(A))\), qui n'est pas en général isomorphe au groupe défini par (1.5.2). Nous changeons donc ici de convention : dans le présent exposé la notation \(k(A)\) désignera exclusivement le groupe défini par (1.5.2).}.
En particulier, pour \(A=C\), on a
\begin{equation}\tag{1.5.3}
 k(D^b(C))=k(C).
\end{equation}

\textbf{Lemme 1.6.} Soient \(C\) une catégorie abélienne, \(A\) une sous-catégorie épaisse (i.e. une sous-catégorie pleine stable par noyaux, conoyaux, extensions). Notons \(D_A(C)\) la sous-catégorie pleine de \(D(C)\) formée des objets \(E\) tels que \(H^i(E)\in\Ob A\) pour tout \(i\). Alors \(D_A(C)\) est une sous-catégorie triangulée \(([\mathrm V]\ \mathrm I\ 1.2.3)\) de \(D(C)\), et le foncteur canonique \(D^b(A)\longrightarrow D_A^b(C)=D^b(C)\cap D_A(C)\) induit un isomorphisme
\[
 k(A)\xrightarrow{\sim} k(D_A^b(C)),
\]
où \(k(A)\) est défini par (1.5.2).

\emph{Preuve.} La première assertion est conséquence immédiate de la suite exacte de cohomologie. Pour la seconde, considérons l'application \(f:\Ob D_A^b(C)\longrightarrow k(A)\) définie par
\[
 f(E)=\sum(-1)^i\cl_A(H^i(E)),
\]
où \(\cl_A\) est l'application (1.5.1). D'après la suite exacte de cohomologie, \(f\) est additive (1.1), donc définit un homomorphisme
\[
 \varphi:k(D_A^b(C))\longrightarrow k(A).
\]
